\section{4 讨论}

\subsection{4.1 主要发现概述}

本研究通过对校园污水管网冬春两季的系统采样分析，共发现47个具有统计学意义的数据模式。
季节比较显示，8个指标存在显著的冬春差异，反映了温度和水文条件对管网碳污染物的深刻影响。
相关分析表明33对变量间的显著关联，为理解碳污染物的多相间迁移机制提供了数据支撑。

\subsection{4.2 碳污染物的季节分异及其影响因素}

CO$\_2$本底值在冬季(638.00)显著高于春季(550.15)(p=0.0001)。
VOCs(ppb)在冬季(598.50)显著高于春季(426.75)(p=0.0003)。
VOCs本底值在冬季(286.70)显著高于春季(141.80)(p=0.0000)。

CO$\_2$与NO$\_2$的正相关可能反应了共同的环境影响因素。温度升高会同时促进有机物分解(产生CO$\_2$)和硝化/反硝化过程(产生NO$\_2$)。此外，管道内微生物活性增强时，有氧呼吸产CO$\_2$和硝化产NO$\_2$同步增加。NO$\_2$作为硝化中间产物，其积累也与DO水平和有机物浓度有关。
Foley et al. (2009)报道管道温室气体排放与微生物活性正相关
有机碳(TOC)是产甲烷菌的主要底物。在厌氧条件下，产甲烷菌通过乙酸发酵或CO$\_2$还原途径将有机碳转化为甲烷。TOC浓度越高，底物越充足，甲烷产生量越大。污水管网中沉积物的有机碳含量是决定甲烷排放的关键因素。
Guisasola et al., 2008, Water Research

上述发现与K. Tang, Z. Yuan（2021）的研究结果一致，进一步验证了该规律的普遍性。

\subsection{4.3 碳氮耦合关系及其环境意义}

总氮（mg/L)与铵态氮（mg/L)呈显著正相关(r=0.985, p=0.0000)。
NaCl(mg/L)与TC(mg/L)呈显著正相关(r=0.657, p=0.0042)。
NaCl(mg/L)与IC(mg/L)呈显著正相关(r=0.709, p=0.0015)。

铵态氮(NH$\_4^+$)与甲烷产生的关联机制包括：(1)氨化作用产生NH$\_4^+$的同时释放有机碳底物；(2)高浓度NH$\_4^+$可抑制产甲烷菌(IC50约3000mg/L)；(3)NH$\_4^+$与CH$\_4$在氮碳耦合中存在协同效应。在低碳氮比条件下，NH$\_4^+$对CH$\_4$的促进作用更明显。
Rajagopal et al., 2013
总氮(TN)与铵态氮(NH$\_4^+$)的强正相关(r=0.985)表明，污水管网中氮的主要存在形式为铵态氮，占TN的大部分。这是因为管道内厌氧条件下，有机氮通过氨化作用转化为NH$\_4^+$，而硝化作用因DO不足而受限，导致NH$\_4^+$积累。TN-NH$\_4^+$的差值主要为有机氮和少量硝态氮。
Guisasola et al. (2008)指出管道厌氧段有机氮氨化是NH$\_4^+$的主要来源

上述发现与Mingjiang Zhang, Chencan Zhang（2025）的研究结果一致，进一步验证了该规律的普遍性。

\subsection{4.4 温室气体排放特征及影响因素}

氧化亚氮(ppm)与NaCl(mg/L)呈显著正相关(r=0.803, p=0.0298)。
CO$\_2$与NO$\_2$（ppm）呈显著正相关(r=0.922, p=0.0000)。
甲烷(ppm)与pH呈显著负相关(r=-0.534, p=0.0224)。

N$\_2$O排放与NaCl浓度的正相关关系可能反应了盐度对硝化过程的影响。高盐度环境下，亚硝酸盐氧化菌(NOB)比氨氧化菌(AOB)更易受抑制，导致亚硝酸盐积累，进而通过硝化反硝化途径产生更多N$\_2$O。此外，NaCl浓度升高会增加污水渗透压，影响微生物群落结构，促进N$\_2$O的产生。
Chandran et al. (2016)报道盐度升高会显著增加污水处理中N$\_2$O排放
DO与硝态氮(NO$\_3^-$)的正相关反应了硝化作用对溶解氧的依赖。硝化细菌(Nitrosomonas和Nitrobacter)是严格好氧菌，DO>2 mg/L时硝化作用活跃，NH$\_4^+$→NO$\_2$-→NO$\_3^-$转化完全。DO<0.5 mg/L时硝化作用几乎停止，NO$\_3^-$浓度极低。因此DO高的采样点NO$\_3^-$浓度也高。
Metcalf \& Eddy (2014)指出DO是硝化作用的限速因子

上述发现与Metcalf \& Eddy, G. Tchobanoglous（2014）的研究结果一致，进一步验证了该规律的普遍性。

\subsection{4.5 盐度对碳氮转化的影响}

氧化亚氮(ppm)与NaCl(mg/L)呈显著正相关(r=0.803, p=0.0298)。
NaCl(mg/L)与电导率(uS/cm)呈显著正相关(r=0.966, p=0.0000)。
NaCl(mg/L)与TC(mg/L)呈显著正相关(r=0.657, p=0.0042)。

N$\_2$O排放与NaCl浓度的正相关关系可能反应了盐度对硝化过程的影响。高盐度环境下，亚硝酸盐氧化菌(NOB)比氨氧化菌(AOB)更易受抑制，导致亚硝酸盐积累，进而通过硝化反硝化途径产生更多N$\_2$O。此外，NaCl浓度升高会增加污水渗透压，影响微生物群落结构，促进N$\_2$O的产生。
Chandran et al. (2016)报道盐度升高会显著增加污水处理中N$\_2$O排放
DO与硝态氮(NO$\_3^-$)的正相关反应了硝化作用对溶解氧的依赖。硝化细菌(Nitrosomonas和Nitrobacter)是严格好氧菌，DO>2 mg/L时硝化作用活跃，NH$\_4^+$→NO$\_2$-→NO$\_3^-$转化完全。DO<0.5 mg/L时硝化作用几乎停止，NO$\_3^-$浓度极低。因此DO高的采样点NO$\_3^-$浓度也高。
Metcalf \& Eddy (2014)指出DO是硝化作用的限速因子

上述发现与Metcalf \& Eddy, G. Tchobanoglous（2014）的研究结果一致，进一步验证了该规律的普遍性。

\subsection{4.6 研究局限与展望}

本研究存在以下局限性：(1)采样时间仅涵盖冬春两季，未纳入夏秋季节数据，难以全面揭示碳污染物的年际变化规律；(2)液相样品的采样点数量有限(n=18)，部分统计分析的统计效力受限；(3)未对管道沉积物和生物膜进行同步采样分析，固相碳的分布特征有待研究。

未来研究可从以下方面拓展：(1)开展四季连续监测，构建碳污染物的完整季节动态模型；(2)结合分子生物学手段(如16S rRNA测序)，揭示驱动碳相间迁移的关键微生物种群；(3)建立管道碳平衡模型，定量评估碳在固-液-气三相之间的转化通量。